\documentclass[12pt]{iptex}
% IDIOMA
% Se necessário, inserir línguas e indicar língua principal (main)
\usepackage[main=brazil,english]{babel}

% Alterando nome do TOC (adicionar em outras línguas, se necessário)
\addto\captionsbrazil{%
	\renewcommand{\contentsname}{SUMÁRIO}
	\renewcommand{\refname}{REFERÊNCIAS}
}
\addto\captionsenglish{%
	\renewcommand{\contentsname}{CONTENTS}
	\renewcommand{\refname}{REFERENCES}
}

% BIBLIOGRAFIA 
\usepackage[abnt-emphasize=bf,alf]{abntex2cite}

% Para bibliografia ABNT com números
%\usepackage[abnt-emphasize=bf,num]{abntex2cite}
%\citebrackets[]

% Para outros estilos o autor deve definir o \bibliographystyle dentro do documento

% Ajustes adicionais
%\setlist[enumerate]{itemsep=3mm} % espaçamento enumerate
%\setlist[itemize]{itemsep=3mm} % espaçamento itemize

% Pacotes adicionais
\usepackage{chemfig}

% Início do documento
\begin{document}

% Parâmetros
%\docNum{}
\tipo{Oficina \LaTeX}
%\cancelaDoc{Relatório Técnico nº 123 456-205} % Descomente para opção cancela e substitui
\cliente{Instituto de Pesquisas Tecnológicas}{IPT}{Av. Professor Almeida Prado, 532}{55555-555}{São Paulo}{SP}
\data{2023}
\titulo{Introdução ao \LaTeX~e aplicação do arquivo de classe IPTeX}
\unidade{Tecnologias Digitais}{TD}
\lab{Seção de Engenharia de Software e Transportes Inteligentes}

% Inserir capa
\capa
\pagestyle{timbrado}

% Inserir resumo
\resumo[Teste de Resumo 1]{Teste de Resumo 2}{Teste de Resumo 3}

\tableofcontents

\vspace{0.5cm}

\pagebreak
\renewcommand{\thepage}{\arabic{page}}
%\setcounter{page}{3}

%\thispagestyle{geral}

\pagestyle{geral}

% Corpo do documento

\section{AGENDA}
\label{sec:agenda}

\begin{itemize}[itemsep=2mm]
	\item \textbf{Apresentação} \textit{(\textcolor{blue}{1° encontro)}}
	\item \textbf{Programas necessários e instalação} \textit{\textcolor{blue}{(1° encontro)}}
	\item \textbf{Estrutura do documento} \textit{\textcolor{blue}{(1° encontro)}}
	\item \textbf{Trabalhando no corpo principal do documento}
	\begin{itemize}[itemsep=2mm]
		\item Seções e subseções \textit{\textcolor{blue}{(2° encontro)}}
		\item Referências internas \textit{\textcolor{blue}{(2° encontro)}}
		\item Listas com marcadores e listas numeradas \textit{\textcolor{blue}{(2° encontro)}}
		\item Floats (tabelas e figuras) \textit{\textcolor{blue}{(3° encontro)}}
		\item Fórmulas \textit{\textcolor{blue}{(4° encontro)}}
		\item Fórmulas químicas \textit{\textcolor{blue}{(4° encontro)}}
	\end{itemize}
	\item \textbf{Utilização do arquivo de classe IPTeX}
	\begin{itemize}[itemsep=2mm]
		\item Tipos de documento implementados \textit{\textcolor{blue}{(5° encontro)}}
		\item Tipos de cabeçalho implementados \textit{\textcolor{blue}{(5° encontro)}}
		\item Parâmetros principais pré definidos \textit{\textcolor{blue}{(5° encontro)}}
		\item Inserção de capa, sumário e resumo, se aplicável \textit{\textcolor{blue}{(6° encontro)}}
		\item Configuração de novos comandos \textit{\textcolor{blue}{(6° encontro)}}
		\item Funções específicas para inserção de tabelas e figuras \textit{\textcolor{blue}{(6° encontro)}}
		\item Assinaturas \textit{\textcolor{blue}{(7° encontro)}}
		\item Anexos \textit{\textcolor{blue}{(7° encontro)}}
		\item Bibliografia \textit{\textcolor{blue}{(7° encontro)}}
		\item Criando novos \textit{templates} \textit{\textcolor{blue}{(8° encontro)}}
	\end{itemize}
	\item \textbf{Possibilidades de automatização com linguagens de programação} \textit{\textcolor{blue}{(8° encontro)}}
\end{itemize}

\clearpage

\section{INTRODUÇÃO}
\label{sec:introducao}

Ao longo dessa oficina serão abordados os principais elementos para a elaboração de documentos a partir de \textit{templates} escritos em \LaTeX.

Será dado enfoque a documentos construídos com o arquivo de classe IPTeX, criado para facilitar a elaboração de documentos técnicos do Instituto de Pesquisas Tecnológicas do Estado de SP.

\section{PROGRAMAS e \textit{LINKS}}
\label{sec:programas}

\subsection{Distribuições e documentação \LaTeX}
\label{subsec:distribuicoes}

Tanto as distribuições quanto a documentação geral podem ser encontradas na página oficial do projeto: \textcolor{blue}{\href{https://www.latex-project.org/}{https://www.latex-project.org/}}.

Para usuários do sistema Windows, são duas opções de distribuição: \textit{MiKTeX} ou \textit{TeX Live}. Na Central de \textit{Software} do IPT pode-se encontrar o \textit{MiKTeX}. Contudo, pode ser mais interessante baixar e instalar o \textit{TeX Live}. O arquivo é bastante grande, mas já vem com grande parte das bibliotecas.

Além da página do projeto, pode-se utilizar o \href{https://en.wikibooks.org/wiki/LaTeX}{\textcolor{blue}{WikiBooks}} para buscar mais informações relativas à documentação.

\subsection{Editor}
\label{subsec:editor}

Existem diversas opções para edição de texto. É possível inclusive usar um bloco de notas, embora não seja a forma mais prática de se trabalhar. Algumas opções são: \textcolor{blue}{\href{https://www.gnu.org/software/emacs/}{\textit{Emacs}}}, \textcolor{blue}{\href{https://code.visualstudio.com/}{\textit{Visual Studio Code}}} e \textcolor{blue}{\href{https://www.texstudio.org/}{\textit{TeXstudio}}}.

Para quem não tem muita familiaridade com \LaTeX, certamente o mais interessante é utilizar o \textit{TeXstudio}. O programa possui interface amigável, com vários assistentes para facilitar o uso de \LaTeX~no dia a dia.

\subsection{Arquivo de classe IPTeX}
\label{subsec:get_iptex}

Atualmente, o IPTeX pode ser obtido diretamente pelo usuário de duas formas: pelo repositório do \textit{GitHub} ou pelo \textit{SharePoint}.

No \textit{GitHub} os arquivos são fechados para membros do IPT (apenas os usuários cadastrados como membros do \href{iptsp}{https://github.com/iptsp}). Basta clonar o projeto IPTeX. Para quem quiser ajudar a desenvolver o projeto, auxiliando na manutenção e nas melhorias necessárias ao longo do tempo, é necessário que tenha uma conta no GitHub com o e-mail do IPT e que faça a solicitação para sua inclusão na organização.

Uma outra forma de acesso, essa aberta para todos os colaboradores do IPT, é por meio do \textit{SharePoint}. Foi criado um site no \textit{SharePoint} ao qual todos têm acesso. Os arquivos nele contidos são os mesmos do projeto no \textit{GitHub}. 

\section{ESTRUTURA}
\label{sec:estrutura}

Vamos considerar o código simples abaixo para entender um pouco melhor a estrutura de um documento escrito em \LaTeX.

\small\begin{verbatim}
	\documentclass[a4paper,12pt]{article} % Classe do documento com opções
	\usepackage[T1]{fontenc} % Pacotes utilizados
	
	% Início do documento
	\begin{document}
		
		Aqui escrevemos o conteúdo do documento

	\end{document}
\end{verbatim}

Nota-se um ambiente delimitado por \verb|\begin{document}....\end{document}|. É nesse local que todos os elementos textuais devem ser alocados (capa, resumo, sumário, seções, bibliografias e anexos). Os elementos anteriores a \verb|\begin{document}| compõem o preâmbulo do documento. É neste lugar em que os arquivos de classe e pacotes são chamados, em que comandos são criados e em que parâmetros de forma são modificados.


\section{SEÇÕES E SUBSEÇÕES}
\label{sec:secoes}

As seções podem ser criadas com os comandos mostrados abaixo. Note que o título do primeiro nível deve ser escrito em caixa alta para que apareça corretamente no Sumário. Além disso, no quarto nível deve-se incluir \verb|\hspace{0pt} \\|, caso contrário não há quebra de linha entre o título da subseção e o texto.

\begin{itemize}
	\item \verb|\section{TÍTULO}| (primeiro nível);
	\item \verb|\subsection{Título}| (segundo nível);
	\item \verb|\subsubsection{Título}| (terceiro nível); e
	\item \verb|\paragraph{Título} \hspace{0pt} \\| (quarto nível).
\end{itemize}


\section{PRIMEIRO NÍVEL}
\label{sec:primeiro_nivel}

\subsection{Segundo nível}
\label{subsec:segundo_nivel}

\subsubsection{Terceiro nível}
\label{subsubsec:terceiro_nivel}

\paragraph{Quarto nível} \hspace{0pt} \\

\section{REFERÊNCIAS INTERNAS}
\label{sec:ref_interna}

As referências internas são feitas com o par \verb|\label{}| (logo após o elemento a ser referenciado), \verb|\ref{}| (ao citar o elemento no texto). Qualquer nome pode ser utilizado entre chaves como referência. Para referenciar o item~\ref{subsubsec:terceiro_nivel} devemos escrever:

\begin{verbatim}
	item~\ref{subsubsec:terceiro_nivel}
\end{verbatim} 

\noindent sendo que o $\sim$ indica que queremos forçar a palavra ``item'' e a referência \ref{subsubsec:terceiro_nivel} na mesma linha.

\section{LISTAS}
\label{sec:listas}

Em \LaTeX~podemos criar listas não ordenadas, com o ambiente \textit{itemize} ou listas ordenadas, com o ambiente \textit{enumerate}. Para mais informações, vide a \href{https://www.overleaf.com/learn/latex/Lists}{\textcolor{blue}{documentação da página do Overleaf}}.

\subsection{Listas não ordenadas}
\label{subsec:itemize}

Uma lista não ordenada é criada com o ambiente\textit{itemize}. Pode-se controlar a separação entre os itens com o parâmetro opcional \textit{itemsep=2mm} (ex. separação de 2 milímetros). Ademais, pode-se controlar todos os marcadores com o parâmetro opcional \textit{label}, ou então individualmente em cada item.

\begin{itemize}[itemsep=1mm, label=$\bullet$]
	\item Primeiro item.
	\item Segundo item.
	\begin{itemize}[itemsep=1mm]
		\item[$\ast$] Subitem dentro do segundo item.
	\end{itemize}
	\item Terceiro item.
	\begin{itemize}[itemsep=1mm]
		\item[--] Subitem dentro do terceiro item.
	\end{itemize}
\end{itemize}

\subsection{\textit{enumerate}}
\label{subsec:enumerate}

Listas ordenadas podem ser criadas no ambiente\textit{enumerate}. O formato padrão é a utilização de algoritmos arábicos. As possibilidades são dadas a seguir.

\begin{itemize}
	\item \verb|\arabic|
	\item \verb|\roman|
	\item \verb|\Roman|
	\item \verb|\Alph|
	\item \verb|\alph|
\end{itemize}

Com o parâmetro opcional \textit{label} pode-se fazer o controle de como os itens são mostrados. O parâmetro \verb|label=\Roman*)|, por exemplo, nos dá um algarismo romano maiúsculo, seguido de parenteses, ou seja,  

\begin{enumerate}[itemsep=1mm, label=\Roman*)]
	\item Primeiro item.
	\item Segundo item.
	\item Terceiro item.
\end{enumerate}


Subitens também podem ser criados da mesma forma que no caso do ambiente \textit{itemize}. Um exemplo é mostrado a seguir.

\begin{enumerate}[itemsep=1mm, label=\Alph*.]
	\item Primeiro item.
	\begin{enumerate}[itemsep=1mm, label=(\roman*)]
		\item Primeiro subitem do primeiro item.
		\item Segundo subitem do primeiro item.
	\end{enumerate}
	\item Segundo item.
	\begin{enumerate}[itemsep=1mm, label=(\roman*)]
		\item Primeiro subitem do segundo item.
		\item Segundo subitem do segundo item.
	\end{enumerate}
\end{enumerate}

\section{FLOATS}
\label{sec:floats}

\textit{Floats} são elementos que não podem ser quebrados entre páginas. Não fazem parte do fluxo normal do texto e ``flutuam'', podendo ser deslocados para a página posterior, por exemplo. Como padrão temos duas classes: figuras (\textit{figure}) e tabelas (\textit{table}). Contudo, novas classes podem ser criadas com o pacote \textit{float}.

Pela característica de ``flutuarem'' no documento, parâmetros posicionais costumam ser utilizados para controlar o posicionamento do elemento de acordo com a necessidade do usuário. Os parâmetros são:

\begin{itemize}
	\item \textbf{t}: na parte superior de uma página de texto;
	\item \textbf{b}: na parte inferior de uma página de texto;
	\item \textbf{h}: na posição do texto em que o \textit{float} aparece (aproximadamente na mesma posição). Caso deseje forçar o elemento exatamente no local inserido, usar \verb|\usepackage{float}| e o parâmetro \textbf{H};
	\item \textbf{p}: em uma página flutuante separada, que não contém texto; e
	\item \textbf{!}: ignora os parâmetros internos que o \LaTeX~usa para determinar o posicionamento do \textit{float}.
\end{itemize}

\subsection{Figuras}
\label{fig:figuras}

A inserção de figuras em \LaTeX~acontece dentro do ambiente \textit{figure}. Figuras são \textit{floats}, ou seja, não podem ser quebradas entre páginas. A Figura~\ref{fig:logo} é um exemplo (se possível, explore o código no arquivo \textit{tex}). 

\begin{figure}[htb]
	\centering
	\captionsetup{justification=raggedright, singlelinecheck=false, width=0.4\textwidth}
	\caption{Logomarca do IPT.}
	\includegraphics[width=0.4\textwidth]{./figuras/headImage.pdf}
	\caption*{Fonte: IPT}
	\label{fig:logo}
\end{figure}


\subsection{Tabelas}
\label{subsec:tabelas}

Um exemplo de tabela usando \LaTeX~é dado abaixo, na Tabela~\ref{tab:exemplo}. Explore os elementos da tabela e veja as alterações provocadas no documento compilado.

\begin{table}[htb!]
	\centering
	\caption{Exemplo de tabela usando \LaTeX.}
	\begin{tabular}{c|c|c}
		\hline
		\textbf{Coluna 1} & \textbf{Coluna 2} & \textbf{Coluna 3} \\
		\hline
		a & b & c \\
		d & e & f \\
		g & h & i \\
		\hline
	\end{tabular}
	\caption*{Fonte: IPT}
	\label{tab:exemplo}
\end{table}

\section{Fórmulas}

\LaTeX~é bastante utilizado na produção de textos e artigos científicos em exatas devido à facilidade de inserção de fórmulas e sua qualidade estética. Abaixo temos dois exemplos da equação de \textit{Navier-Stokes}. A equação~\ref{eq:NS_indicial} mostra a forma indicial (notação de Einstein) para o caso de fluido incompressível com viscosidade constante. Já o sistema de equações~\ref{eq:NS_2D} mostra um caso incompressível, bidimensional, com viscosidade e densidade constantes. 

\begin{equation}
	\rho\frac{Du_i}{Dt} = \rho f_i - \frac{\partial p}{\partial x_i} + \frac{\partial}{\partial x_j}\left[2\mu\left(e_{ij} - \frac{\Delta\delta_{ij}}{3}\right)\right]
	\label{eq:NS_indicial}
\end{equation}

\begin{equation}
\begin{aligned}
	\rho\left(\frac{\partial u}{\partial t} + u\frac{\partial u}{\partial x} + v\frac{\partial u}{\partial y}\right) &= \mu\left[\frac{\partial^2 u}{\partial x^2} + \frac{\partial^2 u}{\partial y^2}\right] - \frac{\partial p}{\partial x} + \rho g_x \\
	\rho\left(\frac{\partial v}{\partial t} + u\frac{\partial v}{\partial x} + v\frac{\partial v}{\partial y}\right) &= \mu\left[\frac{\partial^2 v}{\partial x^2} + \frac{\partial^2 v}{\partial y^2}\right] - \frac{\partial p}{\partial y} + \rho g_y 
\end{aligned}
\label{eq:NS_2D}
\end{equation}

%De acordo com a Equação~\ref{eq:newton2}, $\mathbf{p}$ é a quantidade de movimento. O termo $\upalpha$ e o termo x .... F$_{\text{CDH,ref}}$
%
%\begin{equation}
%	\mathbf{F} = \frac{d\mathbf{p}}{dt} = \frac{d}{dt}(mv)
%	\label{eq:newton2}
%\end{equation}
%
%\begin{equation}
%	y = \alpha x + \beta 
%\end{equation}

\clearpage

\subsection{Fórmulas químicas}
\label{subsec:formulas_quimicas}

Existem alguns pacotes disponíveis. Um deles é o \verb|chemfig|. Alguns exemplos de aplicação são dados abaixo. Para mais exemplos, vide artigo no \href{https://www.overleaf.com/learn/latex/Chemistry_formulae}{\textcolor{blue}{Overleaf}}.

\vspace{0.5cm}

% ângulos absolutos
\chemfig{A-[:50]B-[:-25]C}

\vspace{0.5cm}

% Polígonos regulares
\chemfig{A*5(-B-C-D-E=)}

\vspace{0.5cm}

\chemfig{H-C(-[2]H)(-[6]H)-C(=[1]O)(-[7]H)}

\section{REFERÊNCIAS BIBLIOGRÁFICAS}

Para utilizar referências bibliográficas automáticas deve-se crias um arquivo \textit{BibTeX}, tal qual o \textit{ref.bib} que foi fornecido. Para utilizar ao longo do texto utiliza-se os comandos \verb|\cite{}| ou \verb|\citeonline{}|.

Para facilitar a criação do arquivo \textit{BibTeX} pode-se recorrer ao assistente do \textit{TexStudio} (aba Bibliografia), ao \href{https://scholar.google.com/}{\textcolor{blue}{Google Scholar}} ou a uma ferramenta para gestão de referências bibliográficas, tais quais o \href{https://www.jabref.org/}{\textcolor{blue}{JabRef}}, o \href{https://www.mendeley.com/?interaction_required=true}{\textcolor{blue}{Mendeley}}, o \href{https://www.zotero.org/}{\textcolor{blue}{Zotero}} e o \href{https://www.mybib.com/}{\textcolor{blue}{MyBib}}.

\section{APLICAÇÃO DA CLASSE IPTEX}

A classe IPTeX foi criada para facilitar a utilização do \LaTeX em Relatórios Técnicos. Certificados de Calibração e demais documentos do IPT. O desenvolvimento ocorre de forma contínua no repositório \href{https://github.com/iptsp/IPTeX}{\textcolor{blue}{https://github.com/iptsp/IPTeX}} desde janeiro de 2022.

Para detalhes de utilização veja o manual de instruções disponível no repositório e, caso deseje criar um novo documento técnico, siga um dos modelos criados.

% Bibliografia

%\renewcommand{\refname}{REFERÊNCIAS}
%\addcontentsline{toc}{section}{REFERÊNCIAS}
%\bibliography{ref}

%\label{pag:fim}

\pagebreak
\clearpage

% ANEXO A



\end{document}
